\section{题01：密钥提取（P16783）}

\subsection{从重复数字到代数恒等式}

密钥$A$由$n = 20260606$个数字$9$拼接而成，密钥$B$由$m = 2026$个数字$9$拼接而成。在十进制中，$k$个连续的$9$具有简洁的代数形式：$\underbrace{99\ldots9}_{k\text{个}} = 10^{k} - 1$。取$k=3$可得$999 = 10^3-1$，取$k=5$可得$99999 = 10^5-1$，这一规律对任意$k$均成立。据此，$A = 10^n - 1$，$B = 10^m - 1$，两者的乘积展开为
\[
S = (10^n - 1)(10^m - 1) = 10^{n+m} - 10^n - 10^m + 1.
\]
题目要求提取$S$的十进制表示中（不含前导零，从左往右以$1$为起始编号）第$2026$位、第$20260523$位和第$20262632$位的三个数字$d_1, d_2, d_3$，并以它们组成的三位数$\overline{d_1 d_2 d_3}$作为答案。至此，问题从“两个巨大全$9$数的乘积”转化为了“四个$10$的幂的加减结果中特定位置的数字”，后者可以通过分析代数结构直接求解。

\subsection{暴力乘法的不可行性}

如果试图直接计算$A \times B$，$A$本身约$2 \times 10^7$位十进制数字，存储这个数就需要约$20$ MB内存，乘积$S$更是拥有$n+m = 20262632$位。即便采用基于快速傅里叶变换的大整数乘法，时间复杂度也在$\mathrm{O}(N \log N)$级别，其中$N \approx 2 \times 10^7$为乘积位数，运算量接近$5 \times 10^8$次浮点操作，远超竞赛时限和内存约束。况且这是一道结果填空题，题目设计者显然期望选手绕过数值计算、从数的结构入手。因此放弃暴力乘法、转向代数模式分析是唯一可行的路径。

\subsection{交换减法顺序：一个简洁的推导}

表达式$S = 10^{n+m} - 10^n - 10^m + 1$包含四个项，可以通过不同的分组顺序来化简。一种直观的思路是先合并$10^{n+m} - 10^n$，但该减法涉及跨越多位的借位链，手工追踪较为繁琐。更巧妙的方式是先合并后两项的相反顺序：将$10^{n+m}$与$-10^m$组合。注意到
\[
10^{n+m} - 10^m = 10^m(10^n - 1),
\]
而$10^n - 1$恰好是$n$个$9$，乘以$10^m$相当于在末尾追加$m$个$0$。因此$10^{n+m} - 10^m$的十进制表示极为规则：最高$n$位全部是$9$，紧随其后的$m$位全部是$0$。以小规模参数为例：取$n = 5$，$m = 3$，则$10^3 \times (10^5 - 1) = 10^3 \times 99999 = 99999000$，恰好是$5$个$9$后接$3$个$0$，共$8$位，最高位是$9$，最低位是$0$。

接下来从这个“$n$个$9$、$m$个$0$”的数中减去$10^n$。$10^n$在一个$n+m$位的大数中对齐到$10^n$的幂次位置（从右数第$n+1$位，即从左数第$m$位），该位在上述数中的原值是$9$（因为它落在前$n$位的$9$区域之内——这里依赖$m < n$，本题中$2026 \ll 20260606$显然满足）。由于被减位本就是$9$而非$0$，减法极其简单：$9 - 1 = 8$，且不产生任何向前借位。减完后，第$m$位从$9$变为$8$，其余位不变。最后加上末尾的$+1$：最低位（第$n+m$位）从$0$变为$1$，且因为$n+m-1$位也是$0$，$1$加到$0$上不产生进位。

仍以$n=5, m=3$为例走完整条推导链。第一步，$10^3(10^5-1) = 99999000$（位置$1\sim 5$为$9$，$6\sim 8$为$0$）。第二步，减去$10^5 = 100000$：$10^5$的$1$落在第$3$位（因为总位数$8$，$10^5$对应幂次$5$，从左数位置$8-5=3$），该位置原为$9$，减$1$得$8$；其余位不受影响。结果变为$99899000$。第三步，加$1$：末位$0$变为$1$，得$S = 99899001$。逐位罗列这些数字：$9, 9, 8, 9, 9, 0, 0, 1$。由此可以归纳出$S$在任何一般参数下的完整数字布局。

\subsection{六大区间：积的完整数字模式}

将上述推导推广到一般的$n$和$m$（设$m < n$，本题天然满足），$S$的$n+m$位十进制数字从高位到低位可划分为六个连续区间。第一区间，第$1$位至第$m-1$位，全部是$9$（减法中唯一改变的是第$m$位，这些更高位的$9$完全未受影响）。第二区间，仅含第$m$位，其值为$8$（被$10^n$减走了$1$）。第三区间，第$m+1$位至第$n$位，全部是$9$（同属前$n$个$9$的块，但位于第$m$位右侧，同样未受减法影响）。第四区间，仅含第$n+1$位，其值为$0$（这是后$m$个$0$的开头，$+1$的进位无法到达这里）。第五区间，第$n+2$位至第$n+m-1$位，全部是$0$（同为后$m$个$0$的延续）。第六区间，仅含最末位第$n+m$位，其值为$1$（来自末尾的$+1$）。

为增强信心，用另一组参数交叉验证：取$n = 7$，$m = 4$（总位数$11$）。按相同步骤，$10^4(10^7-1) = 99999990000$，减$10^7$（落在第$4$位）得$99989990000$，加$1$得$S = 99989990001$。观察其$11$位：位置$1\sim 3$为$9$、$4$为$8$、$5\sim 7$为$9$、$8$为$0$、$9\sim 10$为$0$、$11$为$1$，与上述六个区间的划分完全吻合。

\subsection{定位三位目标数字}

将题目参数$n = 20260606$，$m = 2026$代入，总位数$n+m = 20262632$。首先确定第一位$d_1$，位于第$2026$位。该位置恰好等于$m$，落入第二区间（唯一的$8$所在位），故$d_1 = 8$。其次确定第二位$d_2$，位于第$20260523$位。该位置满足$m < 20260523 \le n$（具体而言$20260523 = n - 83$），落入第三区间（$9$的区域），故$d_2 = 9$。最后确定第三位$d_3$，位于第$20262632$位。该位置恰好等于总位数$n+m$，落入第六区间（末尾的$1$），故$d_3 = 1$。因此所求三位整数为$\overline{d_1 d_2 d_3} = \overline{891} = 891$。

\subsection{常数时间的优雅解法}

整个求解过程不涉及任何大整数运算，仅使用整数$n$和$m$的加减与大小比较——将两个全$9$数识别为$10^k - 1$、展开乘积后交换减法顺序以避免借位、归纳六个数字区间、判定三个给定位置各自落入哪个区间。每一步均为$\mathrm{O}(1)$的整数操作，空间开销也仅需存储几个整型变量，同为$\mathrm{O}(1)$。将一个表面上看需要处理约$2 \times 10^7$位数字的“天文”问题化为常数时间，核心在于两条洞察：第一，$k$个$9$恒等于$10^k - 1$，将大数乘法转化为$10$的幂的加减；第二，交换减法顺序（先做$10^{n+m}-10^m$再做$-10^n$而非先做$-10^n$），恰好使被减位落在$9$上从而避免复杂的借位链。这种“用代数结构代替数值计算”的策略，正是竞赛中面对超大数字问题的标准思路。

\subsection{代码实现}
\inputminted{c++}{../src/p01_key_extraction.cpp}
